\section{Detailed Model and Proofs}
\label{sec:detailed_model}

We ask: \emph{Given a local wind disturbance, what is the minimum decision
scope that must be released to recover mission feasibility without
unnecessary replanning?} The short-paper scope includes only wind-induced
perturbations of flight travel time and energy consumption. No-fly zones,
including weather-induced no-fly zones, are excluded from the model.
Recovery is a hierarchical mixed discrete--continuous optimization problem.
Only wind entries change between the paired evaluations; they affect
flight time and energy and thereby subsequent mission constraints.
Disturbance propagation is evaluated on the
fixed incumbent, whereas release sufficiency is an existence statement
about alternative plans in a declared hierarchy.

\subsection{Information, Mission Requirements, and Decision Domain}
\label{subsec:information_remaining}

Let $N,M$ be positive integers. Let $v_0$ denote the depot,
$\mathcal V=\{1,\ldots,N\}$ the relief tasks, and
$\mathcal U=\{1,\ldots,M\}$ the UAVs. The mission distributes one divisible
commodity, with unit mass $m>0$. Intervention levels are $k$,
and decision-component indices are
$a\in\mathcal D_{\mathrm{cmp}}:=\{A,S,P,L\}$.
The fixed mission horizon is $T>0$. Task $i$ has known location
$\mathbf p_i\in\mathbb R^2$ and initial demand $d_i\ge0$;
$\mathbf p_0$ is the depot location. All delivery, loading, onboard,
demand, and stock quantities are expressed in units of the commodity.
The bounded polygonal operating region $\mathcal W\subset\mathbb R^2$ contains all
sites and has diameter $D_{\mathcal W}>0$.

At epoch $t_{n+1}\in[0,T)$, let $H_{n+1}$ contain the executed trajectory
and delivery history, and let $q_i^{\mathrm{del}}(t_{n+1})$ denote the
quantity already delivered. We assume consistent history,
$0\le q_i^{\mathrm{del}}(t_{n+1})\le d_i$, and define
\begin{equation}
d_{i,n+1}^{\mathrm{rem}}
=d_i-q_i^{\mathrm{del}}(t_{n+1}).
\label{eq:residual_demand}
\end{equation}
The residual depot stock is $S_{n+1}^{\mathrm{dep}}\ge0$.
The mission specification supplies total mandatory quantities
$d_i^{\mathrm{mand}}\in[0,d_i]$ and, where applicable, hard service
deadlines $\bar t_i\in[0,T]$. Their residual obligations are
\begin{equation}
\underline d_{i,n+1}
=\left[d_i^{\mathrm{mand}}
-q_i^{\mathrm{del}}(t_{n+1})\right]_+,
\qquad [z]_+:=\max\{z,0\}.
\label{eq:mandatory_residual}
\end{equation}
Let $\mathcal V^{\mathrm{dl}}\subseteq\mathcal V$ be the tasks with hard
deadlines. Define
\begin{align}
\mathcal V_{n+1}^{\mathrm{mand}}
&=\{i:\underline d_{i,n+1}>0\},
\label{eq:mandatory_tasks}\\
\mathcal V_{n+1}^{\mathrm{rem}}
&=\{i:d_{i,n+1}^{\mathrm{rem}}>0\}.
\label{eq:remaining_tasks}
\end{align}
Thus $\mathcal V_{n+1}^{\mathrm{mand}}\subseteq
\mathcal V_{n+1}^{\mathrm{rem}}$. Mandatory amounts and deadlines
are fixed mission inputs and are identical at all intervention levels.
A zero mandatory amount is an explicit absence of that service obligation.
For missions requiring full demand satisfaction, set
$d_i^{\mathrm{mand}}=d_i$. Feasibility is determined by these service
requirements together with the remaining hard constraints.

For the paired evaluation, fix a finite polygonal cell partition
$\mathcal W=\mathop{\dot\bigcup}_{r\in\mathcal R}\mathcal W_r$, where
$\mathcal R=\{1,\ldots,R\}$. Assign each open boundary edge and each vertex
to exactly one adjoining cell. Fix time knots
$t_{n+1}=t_0<t_1<\cdots<t_J=T$, with $R,J\ge1$,
and let $\mathcal J=\{0,\ldots,J-1\}$.
Use $I_j=[t_j,t_{j+1})$ for $j<J-1$ and
$I_{J-1}=[t_{J-1},T]$. Both wind tables cover the entire common domain
$\mathcal W\times[t_{n+1},T]$ and are constant on each cell--time bin:
\begin{equation}
\widehat{\mathbf w}^{\omega}(\mathbf x,s)
=\widehat{\mathbf w}_{rj}^{\omega}\in\mathbb R^2,\qquad
(\mathbf x,s)\in\mathcal W_r\times I_j,\quad \omega\in\{-,+\}.
\label{eq:wind_field}
\end{equation}
The minus table is the previous received wind description restricted to
this common partition and horizon; the plus table is the new description.
The received wind packet is
\begin{equation}
\widehat{\mathcal E}_{n+1}
=\left(t_{n+1},(\widehat{\mathbf w}_{rj}^{+})_
{(r,j)\in\mathcal R\times\mathcal J}\right).
\label{eq:packet}
\end{equation}
The local disturbance has a declared support
$\mathcal B_\Delta\subsetneq\mathcal R\times\mathcal J$:
\begin{equation}
\begin{split}
\widehat{\mathbf w}_{rj}^{+}
&=\widehat{\mathbf w}_{rj}^{-}+\delta\mathbf w_{rj},\\
\delta\mathbf w_{rj}&=\mathbf0\qquad
\text{if }(r,j)\notin\mathcal B_\Delta.
\end{split}
\label{eq:local_wind_update}
\end{equation}
Thus the update is confined to specified cells and time intervals.
For the airspeeds $v_u>0$ defined below, assume a common bound
\begin{equation}
\|\widehat{\mathbf w}_{rj}^{\omega}\|_2\le\bar w<\min_{u\in\mathcal U}v_u,
\qquad (r,j)\in\mathcal R\times\mathcal J,\quad\omega\in\{-,+\}.
\label{eq:wind_bound}
\end{equation}
This bounds the disturbance to a regime in which every direction remains
trackable. Every spatial cell remains traversable; the partition indexes
wind values and imposes no flight exclusion. Operating geometry,
wind-cell boundaries, time bins, demands,
stock, observed state, commitments, and ground-operation rules are fixed in
the comparison and across release levels. The only revised data are wind
vectors. These tables are supplied estimates for the remaining horizon;
certificates are conditional on them.

The observed fleet state is
\begin{equation}
X_{n+1}^{\mathrm{fleet}}
=\left(\mathbf x_{u,n+1},b_{u,n+1},
q_{u,n+1}^{\mathrm{on}},\xi_{u,n+1}\right)_{u\in\mathcal U},
\label{eq:fleet_state}
\end{equation}
where the entries are position, battery energy, onboard commodity quantity, and
execution mode. Elapsed motion and completed operations have already been
settled at $t_{n+1}$. Any unfinished committed operation is represented by a
fixed leading operation with its \emph{residual} duration and transfers;
completed portions are never repeated.
Let $\vartheta$ collect the fixed mission, vehicle, operation, and
representation parameters defined in this section. The information state is
\begin{equation}
\begin{split}
\mathscr I_{n+1}=\bigl(&\vartheta,H_{n+1},X_{n+1}^{\mathrm{fleet}},
\mathbf d_{n+1}^{\mathrm{rem}},S_{n+1}^{\mathrm{dep}},\\
&\bar\Pi_{n+1}^{\mathrm{rem}},\mathcal C_{n+1}^{\mathrm{fix}},
\widehat{\mathcal E}_{0:n+1}\bigr),
\end{split}
\label{eq:information_state}
\end{equation}
where $\mathbf d_{n+1}^{\mathrm{rem}}=(d_{i,n+1}^{\mathrm{rem}})_{i\in\mathcal V}$
collects task demands, $\bar\Pi_{n+1}^{\mathrm{rem}}$ is the incumbent, and
$\mathcal C_{n+1}^{\mathrm{fix}}$ is the commitment mask defined below.
A deterministic non-anticipative policy satisfies
\begin{equation}
\Pi_{n+1}^{\mathrm{rem}}
=\phi_{n+1}(\mathscr I_0,\ldots,\mathscr I_{n+1}).
\label{eq:nonanticipativity}
\end{equation}
A randomized policy may additionally use an internal random seed independent
of future packets. No continuation certificate asserts validity after an
unreceived packet changes the operating conditions.

A remaining plan consists of four decision components:
\begin{equation}
\Pi^{\mathrm{rem}}=(\mathcal A,\mathcal S,\mathcal P,\mathcal L).
\label{eq:remaining_plan}
\end{equation}
At a fixed epoch write $\Pi=\Pi^{\mathrm{rem}}$ and
$\bar\Pi=\bar\Pi_{n+1}^{\mathrm{rem}}$; bars on individual components
always refer to this same incumbent.
The assignment vector $\mathcal A=(a_i)_{i\in\mathcal V_{n+1}^{\mathrm{rem}}}$
has $a_i\in\mathcal U\cup\{\varnothing\}$. Each served task has one
responsible UAV in the remaining plan; repeated partial deliveries by that
UAV are allowed. Cross-UAV splitting of a task within one remaining plan is
excluded by this assignment convention.

To make projections meaningful when routes change, fix a positive integer
slot bound $\bar n_u$ for each UAV before comparing levels.
The sequence component has fixed coordinates
$\mathcal S=(s_{u,h})_{u\in\mathcal U,\,1\le h\le\bar n_u}$ with
$s_{u,h}\in\mathcal V_{n+1}^{\mathrm{rem}}\cup\{v_0,\partial\}$.
The symbol $\partial$ is inactive padding. Active slots form a prefix of
length $n_u(\Pi)\in\{0,\ldots,\bar n_u\}$; if $n_u(\Pi)>0$, its last site is
$v_0$. Intermediate depot visits permit multiple trips. Every active
relief visit to $i$ is assigned to $a_i=u$, and $a_i=u$ requires at least
one such visit. A null assignment therefore means no planned visit.
It is permitted structurally even for a mandatory task; its positive service
lower bound will reject it through hard feasibility.

Let $s_{u,0}=v_{u,n+1}^{\mathrm{cur}}$ be a fixed virtual origin located at
$\mathbf x_{u,n+1}$. The active route is
\begin{equation}
r_u(\Pi)=(v_{u,n+1}^{\mathrm{cur}},s_{u,1},\ldots,s_{u,n_u(\Pi)}).
\label{eq:current_state_route}
\end{equation}
A route with no active slot is structurally permitted only for an idle UAV
already at the depot. A pending delivery at the current relief site is
represented by a leading visit at that site; its incoming path has zero
length. Fixed residual ground operations likewise occupy leading slots or
the initial depot operation as appropriate to the observed mode.

The path component is
$\mathcal P=(P_{u,h})_{u\in\mathcal U,\,1\le h\le\bar n_u}$ on the same
fixed slot indices. An active $P_{u,h}:[0,1]\to\mathcal W$ is a directed
piecewise-linear curve from the preceding site's location to that of
$s_{u,h}$. Curves use normalized arclength parameterization, so adding
collinear vertices does not create a different decision. A zero-length
path is the constant curve at its endpoint; each inactive path is the
constant curve at $\mathbf p_0$. The quantity component is
$\mathcal L=(x_{u,h},y_{u,h})$: delivery occurs at stop $h\ge1$, and
loading at stop $h\ge0$. All inactive quantities are zero. Onboard payload
is induced by these transfers and is not another decision variable.

For each component $a\in\mathcal D_{\mathrm{cmp}}$, let $D_{n+1}^a$ be
its fixed finite set of remaining coordinates, partitioned as
\begin{equation}
D_{n+1}^a=C_{n+1}^a\mathbin{\dot\cup}M_{n+1}^a.
\label{eq:decision_partition}
\end{equation}
Here $C_{n+1}^a$ and $M_{n+1}^a$ are committed and mutable coordinates.
Executed coordinates occur only in $H_{n+1}$.
Write $\mathcal C_{n+1}^{\mathrm{fix}}=(C_{n+1}^a)_a$ and
$\Pi^A=\mathcal A$, $\Pi^S=\mathcal S$, $\Pi^P=\mathcal P$, and
$\Pi^L=\mathcal L$. For coordinate projection $\rho_B$, commitments require
\begin{equation}
\rho_{C_{n+1}^a}(\Pi^a)
=\rho_{C_{n+1}^a}(\bar\Pi_{n+1}^a),
\qquad a\in\mathcal D_{\mathrm{cmp}}.
\label{eq:commitment_projection}
\end{equation}
For example, a first destination may be committed while its unflown path
is mutable. Coordinate labels, padding conventions, and slot bounds are
fixed across all levels. Reassignment and resequencing must respect path
endpoint consistency as well as any fixed coordinates.

The master domain $\mathcal Y_{n+1}^{\mathrm{rem}}$ consists of the
structurally compatible plans just defined, anchored at the observed fleet
state and satisfying \eqref{eq:commitment_projection}. It enforces finite
representation, endpoint and assignment consistency, nonnegative transfer
quantities, and canonical padding, but does not impose time or
resource feasibility. The incumbent is assumed to belong to this domain,
although it may be operationally infeasible after the update.
A commitment override changes the master domain and is a separate safety
action, not level four.

\subsection{Induced Resource and Wind Evaluation}
\label{subsec:resource_wind_evaluation}

Throughout the remaining formulation, the epoch is fixed and omitted from
the residual-data symbols $d_i^{\mathrm{rem}}$, $\underline d_i$, and
$S^{\mathrm{dep}}$.
The following resource recursions apply only to active stops and legs,
with $n_u=n_u(\Pi)$. Let $q_{u,h}^{-}$ and
$q_{u,h}^{+}$ be the scalar onboard quantities immediately before and
after stop $h$. Set
$q_{u,0}^{-}=q_{u,n+1}^{\mathrm{on}}$ and $x_{u,0}=0$.
The stop and flight recursions are
\begin{align}
q_{u,h}^{+}&=q_{u,h}^{-}-x_{u,h}+y_{u,h},
&&0\le h\le n_u,
\label{eq:residual_flow}\\
q_{u,h+1}^{-}&=q_{u,h}^{+},
&&0\le h<n_u.
\label{eq:flight_payload}
\end{align}
Delivery is allowed only at relief visits. Loading is allowed only at
depot stops before a subsequent leg, including stop zero if physically at
the depot. Thus $y_{u,n_u}=0$ and all deliveries at depot stops are zero.
The resource constraints are
\begin{align}
0\le x_{u,h}&\le q_{u,h}^{-},
\label{eq:delivery_availability}\\
\sum_{u}\sum_{h=0}^{n_u-1}y_{u,h}&\le S^{\mathrm{dep}},
\label{eq:residual_inventory}\\
Q_i(\Pi)&:=\sum_{u,h:\,s_{u,h}=i}x_{u,h},
\label{eq:planned_delivery}\\
\underline d_i&\le Q_i(\Pi)\le d_i^{\mathrm{rem}},
\label{eq:mandatory_service}\\
m q_{u,h}^{-}&\le C_u^{\mathrm{pay}},
\qquad
m q_{u,h}^{+}\le C_u^{\mathrm{pay}}.
\label{eq:payload_constraint}
\end{align}
Here $C_u^{\mathrm{pay}}>0$ is payload capacity, and the sum defining
$Q_i$ ranges over active relief visits. An empty sum is zero.
Returned payload remains onboard; it is not credited back to depot stock.
Consequently the inventory bound also prevents intermediate stock
overdrawing, since all loads are nonnegative and no replenishment is modeled.

Each UAV has fixed airspeed $v_u>0$, empty mass $m_u^0>0$, energy
coefficient $\gamma_u>0$, battery capacity $B_u>0$, and reserve
$\beta_u\in[0,B_u]$. The initial battery satisfies
$0\le b_{u,n+1}\le B_u$ and the initial onboard quantities are nonnegative.
Split each finite piecewise-linear path at the fixed wind-cell boundaries,
using the declared ownership rule also for segments lying on a boundary.
Every resulting nonzero segment $e$ lies in one cell $r(e)$ except possibly
at its endpoints, which do not affect the time integral. This refinement
changes no path decision. For its length $\ell_e>0$ and unit direction
$\mathbf t_e$, let $\mathbf t_e^\perp$ be its left normal. Write
$\widehat{\mathbf w}_{r}^{+}(s)=\widehat{\mathbf w}_{rj}^{+}$ for $s\in I_j$,
and define
\begin{equation}
w_{\parallel,e}(s)=\widehat{\mathbf w}_{r(e)}^{+}(s)^{\mathsf T}\mathbf t_e,
\quad
w_{\perp,e}(s)=\widehat{\mathbf w}_{r(e)}^{+}(s)^{\mathsf T}\mathbf t_e^\perp.
\label{eq:wind_components}
\end{equation}
The constant-airspeed track model uses the forward air-velocity branch:
\begin{equation}
s_{u,e}(s)=w_{\parallel,e}(s)
+\sqrt{v_u^2-w_{\perp,e}(s)^2}.
\label{eq:ground_speed}
\end{equation}
The wind bound gives
$\sqrt{v_u^2-w_{\perp,e}(s)^2}>|w_{\parallel,e}(s)|$,
so $s_{u,e}(s)>0$. The ground velocity is $s_{u,e}\mathbf t_e$ and its
air-relative velocity has norm $v_u$; the triangle and reverse-triangle
inequalities then give
$v_u-\bar w\le s_{u,e}(s)\le v_u+\bar w$.
For entry time
$a\in[t_{n+1},T]$, the segment duration is the unique
$\tau\in[0,T-a]$ satisfying
\begin{equation}
\int_a^{a+\tau}s_{u,e}(s)\,\mathrm ds=\ell_e.
\label{eq:travel_relation}
\end{equation}
Such a duration exists exactly when
$\int_a^T s_{u,e}(s)\,\mathrm ds\ge\ell_e$; otherwise the path cannot
complete by the horizon. Strictly positive speed makes the integrated
distance continuous and strictly increasing, establishing uniqueness.
We adopt the linear flight-energy surrogate
\begin{equation}
E_{u,e}^{\mathrm{fly}}
=\gamma_u(m_u^0+m_{u,e}^{\mathrm{pay}})v_u\tau.
\label{eq:flight_energy}
\end{equation}
Here $\gamma_u$ has units of energy per mass per distance and
$m_{u,e}^{\mathrm{pay}}$ is constant onboard mass during the segment.
For a segment of leg $h$, its payload mass is
$m_{u,e}^{\mathrm{pay}}=mq_{u,h-1}^{+}$.
Leg duration $\tau_{u,h}$ and energy $E_{u,h}^{\mathrm{fly}}$ are sums
over its consecutively traversed segments, with each entry time updated
after the preceding segment. A constant path has zero flight duration and
energy. Since speed is constant within each time bin, the integral can be
evaluated by accumulating bin distances and solving linearly in the final
bin. No temporal sampling approximation is part of this definition.
The minus projection uses the same equations with the minus wind table.

Let $t_{u,h}^{\mathrm{arr}}$ and $b_{u,h}^{-}$ be the arrival time and
battery at stop $h$, initialized by
$t_{u,0}^{\mathrm{arr}}=t_{n+1}$ and $b_{u,0}^{-}=b_{u,n+1}$.
A fixed ground-operation protocol specifies service duration
$\delta_{u,h}^{\mathrm{srv}}\ge0$, service energy
$E_{u,h}^{\mathrm{srv}}\ge0$, and depot preparation duration
$\delta_{u,h}^{\mathrm{dep}}\ge0$. These are known deterministic functions
of site, transfers, entering state, and residual execution mode; they
are fixed model data, not independently chosen waiting times.
Service quantities are credited at completion:
\begin{align}
t_{u,h}^{\mathrm{del}}
&=t_{u,h}^{\mathrm{arr}}+\delta_{u,h}^{\mathrm{srv}},
\label{eq:delivery_time}\\
t_{u,h}^{\mathrm{dep}}
&=t_{u,h}^{\mathrm{del}}+\delta_{u,h}^{\mathrm{dep}},
&&h<n_u,
\label{eq:departure_time}\\
t_{u,h+1}^{\mathrm{arr}}
&=t_{u,h}^{\mathrm{dep}}+\tau_{u,h+1},
&&h<n_u.
\label{eq:time_recursion}
\end{align}
The virtual origin has no new service:
$\delta_{u,0}^{\mathrm{srv}}=E_{u,0}^{\mathrm{srv}}=0$.
Depot service duration and energy are zero. At each departure from the depot, the fixed
preparation protocol loads $y_{u,h}$ and restores the battery to $B_u$;
its duration includes the prescribed loading and charging time, or their
remaining time for an already committed operation. Away from the depot,
including an airborne virtual origin, $\delta^{\mathrm{dep}}=0$ and no
energy is replenished. In particular,
\begin{align}
b_{u,h}^{\mathrm{srv}}&=b_{u,h}^{-}-E_{u,h}^{\mathrm{srv}}\ge0,
\label{eq:service_battery}\\
b_{u,h}^{+}
&=\begin{cases}
B_u,&\text{depot stop with }h<n_u,\\
b_{u,h}^{\mathrm{srv}},&\text{otherwise},
\end{cases}
\label{eq:depot_recharge}\\
b_{u,h+1}^{-}&=b_{u,h}^{+}-E_{u,h+1}^{\mathrm{fly}}\ge0,
&&h<n_u.
\label{eq:battery_recursion}
\end{align}
Every depot arrival after flight also satisfies
$b_{u,h}^{-}\ge\beta_u$ before any recharge. Initial depot preparation may
restore a battery below reserve before departure. A zero-slot route has no
preparation and uses its observed battery.
The ground sites are assumed safe for these operations. Depot preparation
has no queue-induced waiting in the propagation recursion. If a finite bay
count $B^{\mathrm{dep}}\in\mathbb N$ is imposed, simultaneous preparation
intervals must satisfy the separate hard constraint
\begin{equation}
\sum_{(u,h)\in\mathcal H^{\mathrm{dep}}(\Pi)}
\mathbf 1_{[t_{u,h}^{\mathrm{arr}},\,t_{u,h}^{\mathrm{dep}})}(t)
\le B^{\mathrm{dep}},\qquad t\in[t_{n+1},T],
\label{eq:depot_bays}
\end{equation}
where $\mathcal H^{\mathrm{dep}}(\Pi)$ contains the nonterminal depot
operations and $\mathbf 1_B$ is the indicator of a set $B$.
The unconstrained-bay model omits \eqref{eq:depot_bays}. A conflict is
infeasible under the given schedule; it does not introduce an implicit
waiting decision.

For UAV $u$ and a path $P$ started at a specified state
$z=(\mathbf x,t,b,q)$, with $t\ge t_{n+1}$, the flight evaluator
either completes by $T$ or establishes a horizon violation. Identical
geometry entered at different times can have different completion times.
Write $\chi_{u,n+1}(P,z)=\chi_u(P,z;\widehat{\mathcal E}_{n+1})$ and define
\begin{equation}
\chi_{u,n+1}(P,z)=
\begin{cases}
\mathsf{complete},&\text{path finishes by }T,\\
\mathsf{late},&\text{otherwise}.
\end{cases}
\label{eq:flight_status}
\end{equation}
A constant path is complete exactly when its start time is at most $T$.
Completion exactly at $T$ is allowed. If the next segment cannot finish
by $T$, its final duration and energy and subsequent flight records are
unavailable; the late status itself is defined. This is a confirmed time
violation, since the wind table covers all times up to $T$.
Define
\begin{equation}
\operatorname{Comp}_u(P,z;\widehat{\mathcal E}_{n+1})
=\mathbf 1\{\chi_{u,n+1}(P,z)=\mathsf{complete}\}.
\label{eq:flight_completion}
\end{equation}
Flight completion does not imply battery sufficiency.
The deterministic nominal evaluation is denoted by
\begin{equation}
\Xi_{n+1}(\Pi)=
\Gamma(\Pi;X_{n+1}^{\mathrm{fleet}},\widehat{\mathcal E}_{n+1},\vartheta).
\label{eq:induced_execution_state}
\end{equation}
It returns the time, payload, battery, and status records induced
above. A ground operation finishing after $T$ likewise establishes a
horizon violation and stops that UAV's remaining evaluation. Unavailable
downstream values are $\bot$, a distinguished symbol outside their
numerical domains. Write $\operatorname{Complete}(\Xi)=1$ when every
required nominal flight and ground operation finishes by $T$ and all
required nominal records are defined. Completeness alone does not imply
resource or service feasibility.

\subsection{Mission Feasibility Under the Updated Wind}
\label{subsec:hard_feasibility}

The ordinary safe terminal set for UAV $u$ is
\begin{equation}
\begin{split}
\mathcal X_u^{\mathrm{safe}}=\{(\mathbf p_0,t,b,q):\;&
0\le t\le T,\ \beta_u\le b\le B_u,\\
&q\ge0,\quad m q\le C_u^{\mathrm{pay}}\}.
\end{split}
\label{eq:safe_terminal_set}
\end{equation}
Landing elsewhere or releasing a commitment belongs to a separate safety
override. At each initial state and post-service relief state, a reserve
certificate is a path $W$ from that state's position to the depot, with no
intermediate delivery or recharge. Its payload is therefore the current
onboard payload. For UAV $u$ at state $z=(\mathbf x,t,b,q)$, define
\begin{equation}
\begin{split}
\mathscr W_{u,n+1}(z)=\{W:\;&
W(0)=\mathbf x,\quad W(1)=\mathbf p_0,\\
&\operatorname{Comp}_u(W,z;\widehat{\mathcal E}_{n+1})=1,\\
&t+\tau_u(W,z)\le T,\\
&E_u(W,z)+\beta_u\le b\},
\end{split}
\label{eq:reserve_witness_set}
\end{equation}
where $W$ uses the same piecewise-linear path convention as a nominal leg,
and $\tau_u(W,z)$ and $E_u(W,z)$ are evaluated for UAV $u$ under the
same received packet. At the assumed-safe depot, the constant path is a
certificate whenever $t\le T$ and $b\ge\beta_u$.
Let $\mathcal Z^{\mathrm{res}}(\Pi)$ contain UAV--state pairs $(u,z)$:
the initial state outside the depot, the first departure state for a UAV
initially at the depot with an active route, and every post-service relief
state. Zero-slot routes are covered by the terminal constraint.
A candidate requires
\begin{equation}
\mathscr W_{u,n+1}(z)\ne\varnothing,
\qquad (u,z)\in\mathcal Z^{\mathrm{res}}(\Pi).
\label{eq:reserve_constraint}
\end{equation}
The existence of a reserve path is a feasibility condition, not a fifth
intervention component. A supplied path certifies it one-sidedly; failure
of one supplied return path does not disprove existence.

Let $z_{u,h}^{\mathrm{dep}}$ contain the stop location,
$t_{u,h}^{\mathrm{dep}}$, $b_{u,h}^{+}$, and $q_{u,h}^{+}$.
At a relief stop its post-service state similarly contains
$t_{u,h}^{\mathrm{del}}$, $b_{u,h}^{\mathrm{srv}}$, and
$q_{u,h}^{-}-x_{u,h}$.
All active flight paths must satisfy
\begin{equation}
\operatorname{Comp}_u(P_{u,h},z_{u,h-1}^{\mathrm{dep}};
\widehat{\mathcal E}_{n+1})=1.
\label{eq:nominal_completion}
\end{equation}
Hard deadlines apply to completed positive deliveries, not merely arrivals:
\begin{equation}
x_{u,h}>0,\ s_{u,h}=i\in\mathcal V^{\mathrm{dl}}
\ \Longrightarrow\ t_{u,h}^{\mathrm{del}}\le\bar t_i.
\label{eq:time_constraints}
\end{equation}
The terminal state $z_u^{\mathrm{term}}$ is the depot-arrival state of the
last active slot, or the observed state for a zero-slot route.

Let $\mathcal Q_{n+1}^{\mathrm{hard}}$ denote the individual resource,
service, deadline, bay (when imposed), reserve, and terminal constraints just
defined, including the pre-recharge depot-arrival reserve bound.
Its predicates are evaluated with all induced quantities substituted.
On plans with incomplete nominal records, a predicate needing an
unavailable value has status undetermined rather than an assumed truth
value. The hard-feasible set is
\begin{equation}
\begin{split}
\mathcal F_{n+1}=\{\Pi\in\mathcal Y_{n+1}^{\mathrm{rem}}:\;&
\operatorname{Complete}(\Xi_{n+1}(\Pi))=1,\\
&\eqref{eq:nominal_completion}\text{ holds},\\
&g(\Pi)=1\quad\forall g\in\mathcal Q_{n+1}^{\mathrm{hard}}\}.
\end{split}
\label{eq:hard_feasible_set}
\end{equation}
Here the terminal predicates include
$z_u^{\mathrm{term}}\in\mathcal X_u^{\mathrm{safe}}$ for every $u$.
We also write $\mathcal F(\mathscr I_{n+1})=\mathcal F_{n+1}$.
A late nominal path or ground operation excludes the plan. Failure to find
a reserve witness alone does not prove that no such witness exists.
Mathematical nonmembership and failure of a particular search are distinct.

Let $\mathcal F_{n+1}^{-}$ be the same feasible set evaluated with the
previous wind table at the same observed state. We assume
$\bar\Pi\in\mathcal F_{n+1}^{-}$, so any newly lost feasibility in this
paired comparison is attributable to the wind change. The received table
need not invalidate the incumbent. Recovery certifies a \emph{complete
remaining mission} under $\widehat{\mathcal E}_{n+1}$; it asserts no
recursive-feasibility guarantee under future wind updates.
Earlier missed obligations remain in the recorded history; only residual
requirements are certified.


\subsection{Typed Disturbance-Propagation Diagnostic}
\label{subsec:propagation_diagnostic}

Write $\widehat{\mathcal E}^{-}$ for the previous packet's wind table
restricted to the common partition and remaining horizon, and
$\widehat{\mathcal E}^{+}=\widehat{\mathcal E}_{n+1}$ for the received table, with
$\omega\in\{-,+\}$. Both projections begin at the \emph{same observed
state at $t_{n+1}$}; the minus projection is a counterfactual evaluation
using the previous wind, not a reconstruction of an earlier physical
state. The incumbent decisions, fixed operation protocol, and commitments
are identical in the two projections.

For UAV $u$, let $\bar e_{u,h}$ denote its incumbent leg at active slot
$h$, and set $\bar n_u^{\mathrm{act}}=n_u(\bar\Pi)$.
For an active route, the reserve-check indices are
\begin{equation}
\begin{split}
\mathcal H_u^{\mathrm{res}}
={}&\{0\}\\
&{}\cup\{h:1\le h\le\bar n_u^{\mathrm{act}},\
\bar s_{u,h}\in\mathcal V_{n+1}^{\mathrm{rem}}\}.
\end{split}
\label{eq:diagnostic_reserve_indices}
\end{equation}
For a zero-slot route, set $\mathcal H_u^{\mathrm{res}}=\varnothing$.
For $h=0$, let $z_{u,0}^{\omega}$ be the observed initial state outside
the depot, or the first departure state if initially at the depot.
For a relief index $h>0$, $z_{u,h}^{\omega}$ is the post-service state
defined above. These are exactly the state pairs required by
$\mathcal Z^{\mathrm{res}}(\bar\Pi)$; unavailable states are $\bot$.
Supply one fixed return geometry $\bar W_{u,h}$ for each
$h\in\mathcal H_u^{\mathrm{res}}$. At other active slots, set
$\bar W_{u,h}=\varnothing$ to indicate an absent reserve-path argument.

Starting from the observed state, process the fixed initial ground
operation once. For an active route the resulting first departure state is
$Z_{u,1}^{\omega}$. Then, for $1\le h\le\bar n_u^{\mathrm{act}}$,
\begin{align}
R_{u,h}^{\omega}
&=\operatorname{EvalLeg}
 (\bar e_{u,h},Z_{u,h}^{\omega};\widehat{\mathcal E}^{\omega},\vartheta,
 \bar W_{u,h}),
\label{eq:two_packet_projection}\\
Z_{u,h+1}^{\omega}
&=\Phi_{u,h}(Z_{u,h}^{\omega},R_{u,h}^{\omega}),
\qquad 1\le h<\bar n_u^{\mathrm{act}},
\label{eq:projection_transition}
\end{align}
where the transition is used only for $h<\bar n_u^{\mathrm{act}}$.
The leg record includes the
following fixed stop evaluation, and $\Phi_{u,h}$ extracts its next
departure state without applying the stop transfers a second time.
The record includes arrival and
service-completion times, flight energy, battery, payload, and completion
and local constraint statuses. Missing numerical values are $\bot$.
A hard resource violation may coexist with computable diagnostic values;
it is recorded rather than silently repaired.

For every $h\in\mathcal H_u^{\mathrm{res}}$, the supplied return geometry
is held fixed between projections. Its battery margin is
\begin{equation}
\mathfrak m_{u,h}^{\omega}
=b(z_{u,h}^{\omega})
-E_u(\bar W_{u,h},z_{u,h}^{\omega};\widehat{\mathcal E}^{\omega})
-\beta_u,
\label{eq:diagnostic_reserve_margin}
\end{equation}
where $b(z)$ extracts a state's battery. The margin is reported only when
the state, completed flight, and energy are defined; otherwise
it is unavailable with an explicit status. These paths are diagnostic
metadata, not extra mutable plan coordinates. Their failure does not
establish failure of the existential reserve condition
\eqref{eq:reserve_constraint}.
The flight-completion status is also reported; a nonnegative battery
margin alone is not a reserve
certificate. The initial reserve record
($h=0$) is associated with the first leg for reporting and query attribution.
It does not change $Z_{u,1}^{\omega}$ or any nominal transition.

There is one change atom $\delta_{rj}$ for each revised wind entry:
\begin{equation}
\mathcal D_{n+1}=\{\delta_{rj}:(r,j)\in\mathcal B_\Delta,\
\widehat{\mathbf w}_{rj}^{-}\ne\widehat{\mathbf w}_{rj}^{+}\}.
\label{eq:wind_atoms}
\end{equation}
Let $\mathcal Q_{u,h}^{\omega}$ contain the wind queries needed by the leg
record, including nominal flight, associated reserve-path evaluation
(including the initial check at $h=1$), and all time bins used in the
distance integral. If a preceding state is unavailable, use a declared
conservative superset of all possible queries for that record. Time-bin
selection and horizon guards must also appear in the dependency graph.

For each atom $\delta$, let $\operatorname{supp}(\delta)$ be the set of
wind queries whose answers the atom can change. Define
\begin{equation}
\operatorname{contact}(\delta,\bar e_{u,h})
=\mathbf 1\!\left\{
\operatorname{supp}(\delta)\cap
(\mathcal Q_{u,h}^{-}\cup\mathcal Q_{u,h}^{+})\ne\varnothing\right\}.
\label{eq:contact}
\end{equation}
Contact tests whether an evaluation queries a changed cell--time bin.
The union of old and new queries is necessary because a changed entry time
can alter the bins traversed. Query supports are conservative diagnostic
objects and do not assert that every contacted record changes.
The relevant atoms are
\begin{equation}
\mathcal D_{n+1}^{\mathrm{rel}}
=\{\delta\in\mathcal D_{n+1}:
\exists u,h,\ \operatorname{contact}(\delta,\bar e_{u,h})=1\}.
\label{eq:relevant_changes}
\end{equation}

Unrolling the fixed-suffix evaluator yields a finite typed directed acyclic
graph
\begin{equation}
\mathcal G_{n+1}
=(\mathcal N_{n+1},\mathcal E_{n+1}^{G},\kappa),
\label{eq:dependency_graph}
\end{equation}
with disjoint vertex classes
\begin{equation}
\mathcal N_{n+1}
=\mathcal N^\Delta\mathbin{\dot\cup}\mathcal N^{\mathrm{wind}}
\mathbin{\dot\cup}\mathcal N^{\mathrm{state}}
\mathbin{\dot\cup}\mathcal N^{\mathrm{eval}}
\mathbin{\dot\cup}\mathcal N^{\mathrm{qty}}
\mathbin{\dot\cup}\mathcal N^{\mathrm{stat}}.
\label{eq:typed_vertices}
\end{equation}
The first class is $\mathcal N^\Delta=\mathcal D_{n+1}$.
Each change atom carries its old or received wind vector in the
corresponding projection; both values are available.
The other five classes contain wind inputs, entering states,
evaluator records, numerical quantities, and categorical statuses.
The map $\kappa$ assigns each vertex its class. Let $\nu^\omega(v)$ be
the value of vertex $v$ in projection $\omega$.
For the parent set $\operatorname{pa}(v)$, each nonroot vertex obeys the
deterministic structural equation
\begin{equation}
\nu(v)=f_v((\nu(w))_{w\in\operatorname{pa}(v)};
\vartheta,\bar\Pi,\bar W).
\label{eq:structural_equation}
\end{equation}
An arc records an argument of that equation, including time-bin selection
and horizon guards. The graph contains the union of dependencies needed by
both evaluations; inactive branches use the same defined unavailable-value
convention. Every wind-query vertex is attributed to at least one
leg record. Outgoing arcs from a change atom lead only to wind-query
vertices in records contacted by that atom, and include every such input
dependency. Atoms outside $\mathcal D_{n+1}^{\mathrm{rel}}$ are therefore
isolated. Only change-atom roots may differ between projections; all other
roots, including unchanged wind entries and observed states, agree.

Within this diagnostic, state-transition arcs join successive legs of the
same UAV. Shared-stock and bay constraints are evaluated separately as
coupling constraints; they do not induce hidden queue delays or cross-UAV
state propagation in $\Gamma$. Thus there is no cross-UAV transition arc
in $\mathcal G_{n+1}$. A shared packet atom may nevertheless contact legs of
several UAVs.

Let the seed set and structural closure be
\begin{align}
\mathcal N_{n+1}^{\mathrm{seed}}
&=\mathcal D_{n+1}^{\mathrm{rel}}
\cup\{v\in\mathcal N^{\mathrm{wind}}:
\exists\delta\in\mathcal D_{n+1}^{\mathrm{rel}},\ (\delta,v)\in
\mathcal E_{n+1}^{G}\},
\label{eq:seed_set}\\
\mathcal N_{n+1}^{\mathrm{cl}}
&=\operatorname{Reach}_{\mathcal G_{n+1}}
  (\mathcal N_{n+1}^{\mathrm{seed}}).
\label{eq:structural_closure}
\end{align}
Reachability includes its starting vertices. It describes structural
dependence and does not assert a realized change.

For every non-change vertex, declare a value space $\mathbb V_v$ and metric
$d_v$ on its available values. Numerical coordinates use a fixed
unit-consistent norm; categorical or canonical symbolic records use the
discrete metric. For a mixed record, use the maximum of coordinate metrics
capped at one, assigning distance one when exactly one coordinate value is
unavailable and zero when both are unavailable. This defines a metric even
for records with partially unavailable fields.
Let $\nu^\omega(v)\in\mathbb V_v\cup\{\bot\}$ denote its projected value.
Let $\varepsilon_v\ge0$ be its reporting tolerance, with
$\varepsilon_v=0$ for categorical values. For two available values write
$\Delta_v=d_v(\nu^-(v),\nu^+(v))$. Define the total predicate
\begin{equation}
\mu(v)=
\begin{cases}
0,&\nu^-(v)=\nu^+(v)=\bot,\\
1,&\text{exactly one value is }\bot,\\
\mathbf 1\{\Delta_v>\varepsilon_v\},
 &\nu^-(v),\nu^+(v)\in\mathbb V_v.
\end{cases}
\label{eq:materiality}
\end{equation}
In particular, no subtraction or norm is applied to $\bot$.
Set
\begin{align}
\mathcal N_{n+1}^{\mathrm{mat}}
&=\{v\in\mathcal N_{n+1}^{\mathrm{cl}}\setminus\mathcal N^\Delta:
\mu(v)=1\},
\label{eq:material_support}\\
\mathcal N_{n+1}^{\mathrm{op}}
&=\mathcal N_{n+1}^{\mathrm{mat}}\setminus\mathcal N^{\mathrm{wind}},
\label{eq:operational_support}\\
\mathcal N_{n+1}^{\mathrm{rep}}
&=\mathcal N_{n+1}^{\mathrm{seed}}\cup\mathcal N_{n+1}^{\mathrm{mat}}.
\label{eq:reportable_support}
\end{align}
Reporting tolerances are applied after reachability and do not prune
intermediate dependencies.

Let $N_{u,h}$ contain the leg's entering-state, evaluation, quantity, and
status vertices, together with its wind-query vertices.
For $h=1$ it also contains the initial reserve-check vertices.
Let $N_{u,h}^{\mathrm{in}}$ be its entering-state subset. Define
\begin{align}
\iota_{u,h}^{\mathrm{dir}}
&=\mathbf 1\{\exists\delta\in\mathcal D_{n+1}^{\mathrm{rel}}:
\operatorname{contact}(\delta,\bar e_{u,h})=1\},
\label{eq:direct_flag}\\
\iota_{u,h}^{\mathrm{prop}}
&=\mathbf 1\{h>1,\ \exists v\in N_{u,h}^{\mathrm{in}}:\mu(v)=1\},
\label{eq:propagated_flag}\\
\iota_{u,h}^{\mathrm{mat}}
&=\mathbf 1\{N_{u,h}\cap\mathcal N_{n+1}^{\mathrm{op}}\ne\varnothing\}.
\label{eq:material_flag}
\end{align}
Direct and propagated flags may both equal one. The affected legs and tasks
are
\begin{align}
\mathcal E_{n+1}^{\mathrm{aff}}
&=\{\bar e_{u,h}:N_{u,h}\cap\mathcal N_{n+1}^{\mathrm{rep}}
\ne\varnothing\},
\label{eq:affected_legs}\\
\mathcal V_{n+1}^{\mathrm{aff}}
&=\mathcal V_{n+1}^{\mathrm{rem}}\cap
\{i:i\text{ is incident to a leg in }\mathcal E_{n+1}^{\mathrm{aff}}\}.
\label{eq:affected_tasks}
\end{align}
For an initial virtual origin physically at a relief site, its task label
is used when identifying incidence.

\medskip
\noindent\textbf{Proposition 1 (outside-closure invariance).}
Assume that table differencing covers every changed wind query,
that the graph obeys the atom-to-query incidence rule above and includes
all data and control dependencies of the two evaluations, and that the
evaluator and non-change roots are identical.
Then
\begin{equation}
v\in\mathcal N_{n+1}\setminus
(\mathcal N_{n+1}^{\mathrm{cl}}\cup\mathcal N^\Delta)
\quad\Longrightarrow\quad \nu^-(v)=\nu^+(v).
\label{eq:outside_closure_invariance}
\end{equation}
\emph{Proof.}
Every change atom with an outgoing arc is relevant and seeded; every
unseeded change atom is isolated. All non-change roots agree.
In topological order, a non-change vertex outside the closure has no parent
inside it and cannot have an unseeded change atom as a parent. Its parents
therefore agree by induction. The same structural equation, including its
branch and unavailable-value rules, gives equal outputs.
\hfill$\square$

\subsection{Task Coupling and Reassignment Neighborhood}
\label{subsec:reassignment_coupling}

The incumbent diagnostic supplies $\mathcal V_{n+1}^{\mathrm{aff}}$.
Reassignment coupling concerns alternative plans throughout
$\mathcal Y_{n+1}^{\mathrm{rem}}$ under the received wind table.

\subsubsection{Constraint instances and coordinate support}
Let $\mathcal Q_{n+1}^{\mathrm{tc}}$ be a finite indexed family whose
conjunction expresses the feasibility conditions in
\eqref{eq:hard_feasible_set}. It contains
$\mathcal Q_{n+1}^{\mathrm{hard}}$ together with active-leg and required
ground-operation completion and definedness checks. Use the fixed slot
bounds to index conditional instances; an inactive instance is satisfied.
The instances retain their resource granularity: shared stock, per-task
service bounds, per-stop resource, deadline and reserve checks, per-UAV
terminal checks, and bay capacity when imposed. The all-time bay inequality
is one predicate on the finite set of preparation intervals. Unrelated
constraints are not merged into a single conjunction vertex.

The tagged mutable-coordinate set is
\begin{equation}
\mathcal M_{n+1}
=\mathop{\dot\bigcup}_{a\in\mathcal D_{\mathrm{cmp}}}
\bigl(\{a\}\times M_{n+1}^a\bigr).
\label{eq:tagged_mutable_coordinates}
\end{equation}
For $g\in\mathcal Q_{n+1}^{\mathrm{tc}}$, define its argument record
$\eta_g(\Pi)$ to contain its activation flag and all plan-dependent
arguments in its predicate, with induced-state recursions substituted.
Include index-selection and branch guards and the definedness status of
each argument. An inactive instance has a canonical inactive record;
unavailable arguments use $\bot$. Fixed inputs belong to
$\mathscr I_{n+1}$. For a reserve constraint, the record includes the
induced state governing the existential set $\mathscr W_{u,n+1}(z)$,
rather than the success of a particular diagnostic return path.
This record and the fixed information determine the status of $g$.

A conservative coordinate support $D_{n+1}(g)\subseteq\mathcal M_{n+1}$
is required to satisfy
\begin{equation}
\begin{split}
\rho_{D_{n+1}(g)}(\Pi)=\rho_{D_{n+1}(g)}(\Pi')
\quad\Longrightarrow\quad\eta_g(\Pi)=\eta_g(\Pi'),\\
\Pi,\Pi'\in\mathcal Y_{n+1}^{\mathrm{rem}}.
\end{split}
\label{eq:constraint_coordinate_support}
\end{equation}
Projection uses the tagged coordinates of all four components.
$\mathcal M_{n+1}$ always satisfies this condition. A conservative
construction takes every mutable leaf of the expanded expression for
$\eta_g$, including all branches, slot-activation guards, and transitive
state dependencies over the common domain. For example, a downstream
battery argument includes earlier payload and path coordinates.
Incumbent traces alone do not capture all permitted routes.

Comparing complete plans avoids inferring independence from an impossible
single-coordinate perturbation in the non-Cartesian master domain.
Argument records retain dependencies even when the incumbent constraint
is slack. Smaller supports require a proof of the same condition; removing
a constraint vertex requires proof that it is satisfied throughout the
common domain.

\subsubsection{Task attribution and the incidence graph}
For UAV $u$, let its potential task participants be
\begin{equation}
\mathcal T_{u,n+1}=\{i\in\mathcal V_{n+1}^{\mathrm{rem}}:
\exists\Pi\in\mathcal Y_{n+1}^{\mathrm{rem}},\ a_i=u\}.
\label{eq:potential_task_participants}
\end{equation}
This includes tasks with committed assignments. Define the attribution
$\Lambda_{n+1}(d)$ of a mutable coordinate by
\begin{equation}
\Lambda_{n+1}(d)=
\begin{cases}
\{i\},&d=(A,i),\\
\mathcal T_{u,n+1},&d\text{ is an }S,P,\text{ or }L
\text{ coordinate of }u.
\end{cases}
\label{eq:coordinate_task_attribution}
\end{equation}
The second case includes depot transfers and conservatively covers changed
neighbors, loads, and later sorties. A declared certified superset of
$\mathcal T_{u,n+1}$ is permissible; the full remaining-task set is safe.
Task-free constraints have no task incidence but remain in the feasibility
test.

Define the task support of $g$ by
\begin{equation}
T_{n+1}(g)=\bigcup_{d\in D_{n+1}(g)}\Lambda_{n+1}(d)
\subseteq\mathcal V_{n+1}^{\mathrm{rem}}.
\label{eq:constraint_task_support}
\end{equation}
The construction rules belong to $\vartheta$; their outputs use only
$\mathscr I_{n+1}$ and are fixed before comparing levels, independently of
the candidate. Different declared conservative constructions can yield
different hierarchies.

The undirected bipartite task--constraint graph is
\begin{equation}
\begin{split}
\mathcal B_{n+1}
=\bigl(&\mathcal V_{n+1}^{\mathrm{rem}}
\mathbin{\dot\cup}\mathcal Q_{n+1}^{\mathrm{tc}},
\mathcal E_{n+1}^{\mathrm{tc}}\bigr),\\
\mathcal E_{n+1}^{\mathrm{tc}}
=\{\{i,g\}:\;&g\in\mathcal Q_{n+1}^{\mathrm{tc}},\ i\in T_{n+1}(g)\}.
\end{split}
\label{eq:incidence_graph}
\end{equation}
Constraint vertices with zero or one incident task do not connect distinct
tasks. The task projection is
\begin{equation}
\begin{split}
\mathcal G_{n+1}^{\mathrm{task}}
&=(\mathcal V_{n+1}^{\mathrm{rem}},\mathcal E_{n+1}^{\mathrm{task}}),\\
\mathcal E_{n+1}^{\mathrm{task}}
&=\{\{i,j\}:i\ne j,\ \exists g\in\mathcal Q_{n+1}^{\mathrm{tc}},\
\{i,j\}\subseteq T_{n+1}(g)\}.
\end{split}
\label{eq:task_projection}
\end{equation}
Thus each multi-task support is a hyperedge represented by a constraint
vertex in $\mathcal B_{n+1}$ and by a clique in its task projection.

\subsubsection{Closure and Level-3 assignment release}
For brevity write $U_\ell$ for the epoch-specific iteration
\begin{equation}
\begin{split}
U_0&=\mathcal V_{n+1}^{\mathrm{aff}},\\
U_{\ell+1}&=U_\ell\cup
\bigcup_{g:\,T_{n+1}(g)\cap U_\ell\ne\varnothing}T_{n+1}(g),\\
U_\infty&=\bigcup_{\ell\ge0}U_\ell.
\end{split}
\label{eq:task_closure_iteration}
\end{equation}
It stabilizes after at most
$|\mathcal V_{n+1}^{\mathrm{rem}}|-|U_0|$ strict enlargements. Equivalently,
\begin{equation}
\begin{split}
U_\infty
&=\operatorname{Reach}_{\mathcal B_{n+1}}(U_0)
\cap\mathcal V_{n+1}^{\mathrm{rem}}\\
&=\operatorname{Reach}_{\mathcal G_{n+1}^{\mathrm{task}}}(U_0).
\end{split}
\label{eq:task_closure_reachability}
\end{equation}
Reachability includes the seeds, including isolated affected tasks.
With $\mathcal V_{n+1}^{\mathrm{mut}}=M_{n+1}^A$, the assignment
neighborhood is
\begin{equation}
\mathcal V_{n+1}^{\mathrm{reassign}}
=U_\infty\cap\mathcal V_{n+1}^{\mathrm{mut}}.
\label{eq:reassignment_neighborhood}
\end{equation}

\medskip
\noindent\textbf{Proposition 2 (least constraint-closed task set).}
$U_\infty$ is the unique least set containing $U_0$ that satisfies
\begin{equation}
T_{n+1}(g)\cap U_\infty\ne\varnothing
\quad\Longrightarrow\quad T_{n+1}(g)\subseteq U_\infty,
\qquad g\in\mathcal Q_{n+1}^{\mathrm{tc}}.
\label{eq:constraint_closed_tasks}
\end{equation}
Consequently no declared constraint support crosses the partition
$(U_\infty,\mathcal V_{n+1}^{\mathrm{rem}}\setminus U_\infty)$.
\emph{Proof.}
The iteration is increasing on a finite task set, so it reaches a fixed
point, which satisfies the implication. If $W\supseteq U_0$ satisfies
the same implication, induction gives $U_\ell\subseteq W$ at every step:
any support meeting $U_\ell$ meets $W$ and is therefore contained in $W$.
Hence $U_\infty\subseteq W$. Alternating task--constraint walks give the
first reachability equality; replacing each two-edge step by its task
projection gives the second. \hfill$\square$

Closure must precede the mutable-assignment intersection. For example,
take four tasks, supports $T(g_1)=\{1,2\}$ and $T(g_2)=\{2,3\}$, seed
$U_0=\{1\}$, and mutable assignments $M^A=\{1,3,4\}$.
Then $U_\infty=\{1,2,3\}$ and the reassignment set is $\{1,3\}$.
Deleting task 2 would incorrectly disconnect task 3: its committed
assignment does not fix its route, path, or quantity dependencies.

Enlarging the seeds or any support cannot shrink the closure. Empty seeds
give an empty neighborhood. A reached support containing all mutable tasks
makes the neighborhood equal to $M_{n+1}^A$. Shared stock or bay capacity
can cause this case. Under the stated attribution,
$\mathcal T_{u,n+1}=\mathcal V_{n+1}^{\mathrm{rem}}$ also makes any
constraint depending on a mutable route coordinate of $u$ incident to all
tasks. Local wind therefore does not guarantee local reassignment.

Minimality holds only among sets closed under the declared supports, not
among assignment sets sufficient for repair. Membership permits a move;
it does not require one. The no-crossing property concerns $U_\infty$
before commitment filtering. It does not imply independent subproblems:
structural rules remain in the common domain, and all mutable sequence,
path, and quantity coordinates are released at levels two and three.
Every candidate still requires full constraint validation.

Given the supports, a multi-source traversal of $\mathcal B_{n+1}$ visits
each constraint vertex once, using time and storage
$O(|\mathcal V_{n+1}^{\mathrm{rem}}|+|\mathcal Q_{n+1}^{\mathrm{tc}}|
+\sum_g|T_{n+1}(g)|)$. This excludes support construction and avoids
materializing the task projection.

\subsection{Cumulative Intervention Spaces}
\label{subsec:intervention_spaces}

For $a\in\mathcal D_{\mathrm{cmp}}$, let
$R_{k,n+1}^a\subseteq M_{n+1}^a$ be the released coordinates at level
$k\in\mathbb K:=\{0,1,2,3,4\}$. Suppressing the epoch only inside the
following table, the release schedule is
\begin{equation}
\begin{array}{@{}ccccc@{}}
\toprule
k&R_k^A&R_k^S&R_k^P&R_k^L\\
\midrule
0&\varnothing&\varnothing&\varnothing&\varnothing\\
1&\varnothing&\varnothing&M^P&\varnothing\\
2&\varnothing&M^S&M^P&M^L\\
3&\mathcal V^{\mathrm{reassign}}&M^S&M^P&M^L\\
4&M^A&M^S&M^P&M^L\\
\bottomrule
\end{array}.
\label{eq:release_schedule}
\end{equation}
Let $B_{k,n+1}^a=M_{n+1}^a\setminus R_{k,n+1}^a$. All five spaces have
the common definition
\begin{equation}
\begin{split}
\mathcal X_{k,n+1}=\{\Pi\in\mathcal Y_{n+1}^{\mathrm{rem}}:\;&
\rho_{B_{k,n+1}^a}(\Pi^a)
=\rho_{B_{k,n+1}^a}(\bar\Pi_{n+1}^a),\\
&a\in\mathcal D_{\mathrm{cmp}}\}.
\end{split}
\label{eq:intervention_space}
\end{equation}
Thus $\mathcal X_{0,n+1}=\{\bar\Pi_{n+1}^{\mathrm{rem}}\}$ and
$\mathcal X_{4,n+1}=\mathcal Y_{n+1}^{\mathrm{rem}}$.
Level one changes only uncommitted paths. Level two also changes service
order and loading/delivery quantities, preserving assignment.
Quantities are released there because resequencing changes the route on
which payload is carried and replenished. Level three additionally
releases assignments in the declared coupling closure; level four releases
all mutable assignments. No level changes observed history, current state,
slot bounds, fixed operation rules, or commitments.
In particular, the graph-defined Level-3 space is exactly
\begin{equation}
\begin{split}
\mathcal X_{3,n+1}=\{\Pi\in\mathcal Y_{n+1}^{\mathrm{rem}}:\;&
a_i=\bar a_i\\
&\forall i\in M_{n+1}^A\setminus
\mathcal V_{n+1}^{\mathrm{reassign}}\}.
\end{split}
\label{eq:level_three_space}
\end{equation}
Assignments in the neighborhood may also retain their incumbent values.
Feasible Level-3 repairs belong to
$\mathcal X_{3,n+1}\cap\mathcal F_{n+1}$; graph membership alone gives no
feasibility certificate.

\medskip
\noindent\textbf{Proposition 3 (nesting).}
If all levels share the master domain, incumbent, coordinate convention,
and commitment mask, and released coordinates are cumulative, then
\begin{equation}
\mathcal X_{k,n+1}\subseteq\mathcal X_{k+1,n+1},
\qquad k=0,1,2,3.
\label{eq:nesting}
\end{equation}
\emph{Proof.}
Cumulative release gives
$B_{k+1,n+1}^a\subseteq B_{k,n+1}^a$ for every component.
Every equality required at level $k+1$ is therefore already satisfied at
level $k$; the common domain imposes identical remaining restrictions.
\hfill$\square$

Equality is allowed. In particular, levels three and four coincide whenever
$\mathcal V_{n+1}^{\mathrm{reassign}}=\mathcal V_{n+1}^{\mathrm{mut}}$.
With an empty reassignment neighborhood, levels two and three coincide.
The hierarchy is over permitted changes, not over the number of coordinates
that an eventual candidate actually changes.


\subsection{Minimum Decision Release and Certificates}
\label{subsec:minimum_escalation}

Define $\mathcal F_{k,n+1}:=\mathcal X_{k,n+1}\cap\mathcal F_{n+1}$.
The minimum sufficient release level is
\begin{equation}
k_{n+1}^F=\min\{k\in\mathbb K:\mathcal F_{k,n+1}\ne\varnothing\},
\qquad \min\varnothing:=+\infty.
\label{eq:feasibility_escalation}
\end{equation}
Nesting gives
\begin{equation}
\mathcal F_{k,n+1}\ne\varnothing
\ \Longrightarrow\
\mathcal F_{j,n+1}\ne\varnothing,\qquad j\in\mathbb K,\ j\ge k.
\label{eq:feasibility_monotonicity}
\end{equation}
Consequently, for a finite $k$,
\begin{equation}
\begin{split}
k_{n+1}^F=k\quad\Longleftrightarrow\quad&
\mathcal F_{k,n+1}\ne\varnothing,\\
&\mathcal F_{j,n+1}=\varnothing
\quad\text{for all }j\in\mathbb K,\ j<k.
\end{split}
\label{eq:minimumharacterization}
\end{equation}
This minimum is over the declared five-level hierarchy. Unreleased
coordinates remain equal to the incumbent; released coordinates may also
remain unchanged. The formulation does not claim to find a smallest
arbitrary subset of individual coordinates. Disturbance severity or
diagnostic reachability depth alone does not determine the minimum.

The mathematical absence of a mission-feasible continuation is
\begin{equation}
k_{n+1}^F=+\infty
\quad\Longleftrightarrow\quad
\mathcal F_{4,n+1}=\varnothing.
\label{eq:no_feasible_continuation}
\end{equation}
A computational procedure may report no feasible continuation only with
a valid certificate of this emptiness; its status code is
\begin{equation*}
\mathsf{NO\_FEASIBLE\_CONTINUATION}.
\end{equation*}
Such emptiness is relative to the stated representation,
commitments, mandatory service, ground-operation protocol, and received
wind table. It does not establish impossibility after future wind
information arrives or after the domain is changed by an emergency override.

A supplied plan $\widetilde\Pi_k$, together with feasible reserve-path
witnesses and its induced-state checks, can establish
\begin{equation}
\widetilde\Pi_k\in\mathcal F_{k,n+1}
\quad\Longrightarrow\quad k_{n+1}^F\le k.
\label{eq:candidate_upper_bound}
\end{equation}
The converse is false for a particular candidate. Failed construction,
rejected validation, and a time or iteration limit are not certificates
that $\mathcal F_{k,n+1}$ is empty.

More generally, let $\mathbb K^{\mathrm{wit}}$ contain the levels with
validated feasible witnesses, and $\mathbb K^{\mathrm{emp}}$ those with
valid emptiness certificates. Define
\begin{align}
u_{n+1}^F&=\min\mathbb K^{\mathrm{wit}},
\label{eq:certified_upper}\\
\ell_{n+1}^F
&=\begin{cases}
+\infty,&4\in\mathbb K^{\mathrm{emp}},\\
1+\max(\{-1\}\cup\mathbb K^{\mathrm{emp}}),&\text{otherwise}.
\end{cases}
\label{eq:certified_lower}
\end{align}
The empty minimum convention also applies to $u_{n+1}^F$.
Then
\begin{equation}
\ell_{n+1}^F\le k_{n+1}^F\le u_{n+1}^F.
\label{eq:certificate_interval}
\end{equation}
\emph{Justification.}
A witness proves nonemptiness at its level. If level $j$ is empty,
nesting makes every lower level empty; if level four is empty, all levels
are empty. These two implications give the bounds. Equality of finite
bounds certifies the minimum. An empty witness set alone gives only the
vacuous upper bound $+\infty$. \hfill$\square$

The statuses \emph{witness found}, \emph{space certified empty}, and
\emph{search unresolved} concern a search problem. They are distinct from
the flight statuses complete and late in
\eqref{eq:flight_status}.
Likewise, an emptiness proof over a restricted path library is not an
emptiness proof over the entire piecewise-linear path domain. Its
representation restriction must accompany the certificate.

\balance
\subsection{Optional Intervention Magnitude}
\label{subsec:intervention_magnitude}

When $k_{n+1}^F<+\infty$, every plan in
$\mathcal F_{k_{n+1}^F,n+1}$ is an acceptable restored plan.
The distance below is optional for reporting or selecting among these
witnesses. It is not a feasibility constraint and does not alter
$k_{n+1}^F$.

Write $\bar\Pi=\bar\Pi_{n+1}^{\mathrm{rem}}$ and
$H_{\mathrm{slot}}=\sum_u\bar n_u>0$. Define
$N_{\mathrm{rem}}=|\mathcal V_{n+1}^{\mathrm{rem}}|$ and the fixed quantity
scale $\sigma=1+\sum_i d_i$, expressed in commodity units.
Write
\begin{equation*}
\Delta_{u,h}^{P}=\sup_{\theta\in[0,1]}
\|P_{u,h}(\theta)-\bar P_{u,h}(\theta)\|_2.
\end{equation*}
The common coordinate representation gives the explicit component distances
\begin{align}
d_A&=\frac{\sum_{i\in\mathcal V_{n+1}^{\mathrm{rem}}}
\mathbf 1\{a_i\ne\bar a_i\}}{\max\{1,N_{\mathrm{rem}}\}},
\label{eq:assignment_distance}\\
d_S&=\frac{1}{H_{\mathrm{slot}}}
\sum_u\sum_{h=1}^{\bar n_u}\mathbf 1\{s_{u,h}\ne\bar s_{u,h}\},
\label{eq:sequence_distance}\\
d_P&=\frac{1}{H_{\mathrm{slot}}D_{\mathcal W}}
\sum_u\sum_{h=1}^{\bar n_u}
\Delta_{u,h}^{P},
\label{eq:path_distance}\\
d_L&=\frac{\displaystyle\sum_u\sum_{h=0}^{\bar n_u}
\bigl(|x_{u,h}-\bar x_{u,h}|+|y_{u,h}-\bar y_{u,h}|\bigr)}
{2\sigma(H_{\mathrm{slot}}+M)}.
\label{eq:quantity_distance}
\end{align}
Here $x_{u,0}=0$ in both plans and all inactive entries use their
canonical padding. Sequence distance measures slot changes.
Path distance compares geometry under the fixed arclength convention,
including its endpoints; it is not a binary path-change flag.
For fixed positive weights $\alpha_a$ satisfying $\sum_a\alpha_a=1$, set
\begin{equation}
d_{n+1}(\Pi,\bar\Pi)=\alpha_A d_A+\alpha_S d_S+\alpha_P d_P+\alpha_L d_L.
\label{eq:intervention_distance}
\end{equation}
It is finite and nonnegative, and vanishes exactly when the represented
remaining decisions agree. All weights and scales are held fixed across
levels. Neither alternative waypoint encodings of the same arclength
curve nor inactive padding create spurious path changes.

If $k_{n+1}^F<+\infty$, let
\begin{equation}
d_{n+1}^{F,\star}
=\inf_{\Pi\in\mathcal F_{k_{n+1}^F,n+1}}d_{n+1}(\Pi,\bar\Pi).
\label{eq:minimum_disruption}
\end{equation}
When the infimum is attained, a least-disruptive continuation belongs to
the corresponding $\argmin$. Without an attainment assumption, for any
$\delta_d>0$ one may select
\begin{equation}
\Pi_{n+1}^{F,\delta_d}\in\mathcal F_{k_{n+1}^F,n+1},
\qquad
d_{n+1}(\Pi_{n+1}^{F,\delta_d},\bar\Pi)
\le d_{n+1}^{F,\star}+\delta_d.
\label{eq:secondary_selection}
\end{equation}
Such a plan exists by the definition of infimum. Scope remains the primary
criterion; distance cannot trade a broader level for a smaller change.

The formulation therefore separates diagnosis on the fixed incumbent from
existence of a mission-feasible continuation. A valid witness establishes
feasibility at its declared level; determining the minimum also requires
certified exclusion of lower levels. These statements are relative to the
received information, commitments, and declared representation.
